\documentclass[11pt,a4paper]{article}

% Packages
\usepackage[utf8]{inputenc}
\usepackage[margin=1in]{geometry}
\usepackage{hyperref}
\usepackage{listings}
\usepackage{xcolor}
\usepackage{graphicx}
\usepackage{fancyhdr}
\usepackage{tocloft}
\usepackage{titlesec}

% Code listing style
\lstset{
    basicstyle=\ttfamily\small,
    keywordstyle=\color{blue}\bfseries,
    commentstyle=\color{green!50!black},
    stringstyle=\color{red},
    showstringspaces=false,
    breaklines=true,
    frame=single,
    numbers=left,
    numberstyle=\tiny\color{gray},
    language=C++,
    tabsize=4
}

% Header and footer
\pagestyle{fancy}
\fancyhf{}
\lhead{CGI Implementation Documentation}
\rhead{Webserv Project}
\cfoot{\thepage}

% Title
\title{\textbf{CGI Implementation Documentation}\\
\large RFC 3875 Compliant Common Gateway Interface Handler\\
\vspace{0.5cm}
\normalsize Webserv - 42 School Project}
\author{QuantumBytes Team}
\date{\today}

\begin{document}

\maketitle
\newpage

\tableofcontents
\newpage

\section{Introduction}

\subsection{Overview}

This document provides comprehensive documentation for the CGI (Common Gateway Interface) implementation in the Webserv project. The implementation is designed to execute external scripts and programs in response to HTTP requests, following the RFC 3875 specification.

\subsection{Purpose and Scope}

The CGI handler enables the web server to:

\begin{itemize}
    \item Execute external scripts (Python, PHP, Perl, shell scripts, etc.)
    \item Pass HTTP request information via environment variables
    \item Handle bidirectional communication with CGI processes
    \item Provide secure script execution with validation and sanitization
    \item Support both GET and POST HTTP methods
\end{itemize}

\subsection{RFC 3875 Compliance}

This implementation adheres to RFC 3875 - "The Common Gateway Interface (CGI) Version 1.1" with the following features:

\begin{itemize}
    \item \textbf{Environment Variables}: All required CGI meta-variables are implemented
    \item \textbf{Request Parsing}: Proper URI parsing including SCRIPT\_NAME, PATH\_INFO, and QUERY\_STRING
    \item \textbf{Header Handling}: HTTP headers converted to HTTP\_* environment variables
    \item \textbf{Content Handling}: Proper CONTENT\_TYPE and CONTENT\_LENGTH processing
    \item \textbf{Output Parsing}: RFC-compliant header/body separation
\end{itemize}

\subsection{42 School Project Context}

This implementation is part of the 42 School Webserv project with specific requirements:

\begin{itemize}
    \item \textbf{C++98 Standard}: No C++11 or later features
    \item \textbf{No External Libraries}: Pure C++98 with standard library only
    \item \textbf{Error Handling}: Robust system call error checking
    \item \textbf{Non-blocking I/O}: Pipes configured for non-blocking operation
    \item \textbf{No errno in Production}: errno usage limited to test code only
\end{itemize}

\subsection{Project Structure}

The implementation spans multiple components:

\begin{itemize}
    \item \texttt{dependencies/libhttp/} - HTTP library containing CGI class
    \item \texttt{dependencies/libutils/} - Reusable utility functions
    \item \texttt{tests/} - Comprehensive test suite (50 tests)
    \item \texttt{tests/cgi-bin/} - Sample CGI scripts for testing
\end{itemize}

\newpage
\section{Architecture and Design}

\subsection{Class Structure}

The CGI implementation is encapsulated in a single \texttt{CGI} class with clear separation of concerns:

\begin{lstlisting}[caption=CGI Class Overview]
class CGI {
public:
    enum SetupResult { /* 7 error codes */ };
    
    CGI();
    ~CGI();
    
    SetupResult setup(...);
    pid_t execute(...);
    static bool parseCGIOutput(...);
    static bool isCGIScript(...);
    
private:
    // State
    std::map<std::string, std::string> _env;
    std::vector<char*> _envp;
    std::string _scriptPath;
    std::string _interpreterPath;
    std::string _requestBody;
    
    // Validation
    SetupResult validateInputs(...);
    
    // Environment building
    void buildEnvironment(...);
    void addEnvironmentVariable(...);
    
    // Pipe management
    bool createPipes(...);
    void setupChildProcess(...);
};
\end{lstlisting}

\subsection{Library Separation}

The project follows DRY (Don't Repeat Yourself) principle by separating utilities:

\subsubsection{libutils (General Utilities)}

Located in \texttt{dependencies/libutils/}, provides reusable functions in the \texttt{utils::} namespace:

\begin{itemize}
    \item \textbf{String Utilities}: trim(), toUpper(), toLower()
    \item \textbf{Path Utilities}: getDirectory(), getFilename(), getExtension()
    \item \textbf{Validation}: isValidPort(), containsPathTraversal(), sanitizeForEnv()
    \item \textbf{Conversion}: intToString(), sizeToString(), stringToInt()
\end{itemize}

\subsubsection{libhttp (HTTP-Specific)}

Located in \texttt{dependencies/libhttp/}, contains HTTP-related classes:

\begin{itemize}
    \item \texttt{CGI} - CGI execution handler
    \item \texttt{Request} - HTTP request parser
    \item \texttt{Response} - HTTP response builder
    \item \texttt{Utils} - HTTP-specific utilities (methodToString())
\end{itemize}

\subsection{Design Principles}

\subsubsection{Single Responsibility Principle (SRP)}

Each method has a single, well-defined responsibility:

\begin{itemize}
    \item \texttt{setup()} - Validate inputs and prepare environment
    \item \texttt{execute()} - Fork and execute CGI script
    \item \texttt{parseCGIOutput()} - Parse CGI output headers
    \item \texttt{validateInputs()} - Input validation only
    \item \texttt{buildEnvironment()} - Environment variable construction only
    \item \texttt{createPipes()} - Pipe creation and setup only
\end{itemize}

\subsubsection{Keep It Simple, Stupid (KISS)}

\begin{itemize}
    \item Simple, focused interface with 3 main public methods
    \item Clear error reporting via SetupResult enum
    \item Straightforward pipe management
    \item Direct system calls without unnecessary abstractions
\end{itemize}

\subsubsection{Fail Fast}

\begin{itemize}
    \item All validation occurs in \texttt{setup()} before execution
    \item Early return on validation errors
    \item Clear error messages for each failure case
    \item No execution attempt if setup fails
\end{itemize}

\newpage
\section{API Reference}

\subsection{SetupResult Enumeration}

Error codes returned by the \texttt{setup()} method:

\begin{lstlisting}[caption=SetupResult Enum]
enum SetupResult {
    SETUP_OK = 0,
    SETUP_ERROR_EMPTY_SCRIPT_PATH,
    SETUP_ERROR_INVALID_PORT,
    SETUP_ERROR_PATH_TRAVERSAL,
    SETUP_ERROR_SCRIPT_NOT_FOUND,
    SETUP_ERROR_SCRIPT_NOT_EXECUTABLE,
    SETUP_ERROR_INTERPRETER_NOT_FOUND
};
\end{lstlisting}

\begin{itemize}
    \item \textbf{SETUP\_OK}: Validation successful, ready to execute
    \item \textbf{SETUP\_ERROR\_EMPTY\_SCRIPT\_PATH}: Script path is empty
    \item \textbf{SETUP\_ERROR\_INVALID\_PORT}: Port number outside 1-65535 range
    \item \textbf{SETUP\_ERROR\_PATH\_TRAVERSAL}: Path contains ".." sequences or null bytes
    \item \textbf{SETUP\_ERROR\_SCRIPT\_NOT\_FOUND}: Script file does not exist
    \item \textbf{SETUP\_ERROR\_SCRIPT\_NOT\_EXECUTABLE}: Script lacks execute permission
    \item \textbf{SETUP\_ERROR\_INTERPRETER\_NOT\_FOUND}: Interpreter not found or not executable
\end{itemize}

\subsection{Public Methods}

\subsubsection{setup()}

Validates inputs and prepares CGI execution environment.

\begin{lstlisting}[caption=setup() Method]
SetupResult setup(
    const libhttp::Request &request,
    const std::string &scriptPath,
    const std::string &interpreterPath,
    const std::string &serverName,
    int serverPort
);
\end{lstlisting}

\textbf{Parameters:}
\begin{itemize}
    \item \texttt{request} - HTTP request object containing method, URI, headers, and body
    \item \texttt{scriptPath} - Absolute path to CGI script file
    \item \texttt{interpreterPath} - Path to interpreter (e.g., /usr/bin/python3), or empty for direct execution
    \item \texttt{serverName} - Server hostname for SERVER\_NAME variable
    \item \texttt{serverPort} - Server port number (1-65535)
\end{itemize}

\textbf{Return Value:}
\begin{itemize}
    \item \texttt{SETUP\_OK} on success
    \item Specific error code on failure (see SetupResult enum)
\end{itemize}

\textbf{Behavior:}
\begin{enumerate}
    \item Validates all inputs (path, port, file existence, permissions)
    \item Builds CGI environment variables from request
    \item Prepares environment array for execve()
    \item Sets internal state to ready for execution
\end{enumerate}

\subsubsection{execute()}

Forks and executes the CGI script.

\begin{lstlisting}[caption=execute() Method]
pid_t execute(int &readPipeFd, int &writePipeFd);
\end{lstlisting}

\textbf{Parameters:}
\begin{itemize}
    \item \texttt{readPipeFd} - Output: file descriptor to read CGI output
    \item \texttt{writePipeFd} - Output: file descriptor to write request body
\end{itemize}

\textbf{Return Value:}
\begin{itemize}
    \item PID of child process on success
    \item -1 on error (check getError() for details)
\end{itemize}

\textbf{Critical Usage Notes:}
\begin{enumerate}
    \item Must call \texttt{setup()} successfully before \texttt{execute()}
    \item Caller MUST close \texttt{writePipeFd} after writing request body
    \item For GET requests, close \texttt{writePipeFd} immediately
    \item For POST requests, write body then immediately close \texttt{writePipeFd}
    \item Failure to close will cause CGI script to hang waiting for EOF
\end{enumerate}

\subsubsection{parseCGIOutput()}

Parses CGI output into headers and body.

\begin{lstlisting}[caption=parseCGIOutput() Static Method]
static bool parseCGIOutput(
    const std::string &output,
    std::map<std::string, std::string> &headers,
    std::string &body
);
\end{lstlisting}

\textbf{Parameters:}
\begin{itemize}
    \item \texttt{output} - Raw CGI output from script
    \item \texttt{headers} - Output: parsed HTTP headers
    \item \texttt{body} - Output: response body
\end{itemize}

\textbf{Return Value:}
\begin{itemize}
    \item \texttt{true} if headers parsed successfully
    \item \texttt{false} if no header separator found (treats entire output as body)
\end{itemize}

\textbf{Header Format:}
\begin{itemize}
    \item Headers separated by blank line (CRLF or LF)
    \item Format: "Header-Name: Value"
    \item Both key and value are trimmed for liberal parsing
\end{itemize}

\subsubsection{isCGIScript()}

Checks if a file is a CGI script based on extension.

\begin{lstlisting}[caption=isCGIScript() Static Method]
static bool isCGIScript(
    const std::string &path,
    const std::string &extension
);
\end{lstlisting}

\textbf{Parameters:}
\begin{itemize}
    \item \texttt{path} - File path to check
    \item \texttt{extension} - Extension to match (e.g., ".cgi", ".py", ".pl")
\end{itemize}

\textbf{Return Value:}
\begin{itemize}
    \item \texttt{true} if path ends with the specified extension
    \item \texttt{false} otherwise
\end{itemize}

\subsubsection{getError()}

Returns error message from last operation.

\begin{lstlisting}[caption=getError() Method]
const std::string &getError() const;
\end{lstlisting}

\textbf{Return Value:}
\begin{itemize}
    \item Error message string (empty if no error)
\end{itemize}

\subsubsection{getSetupResultMessage()}

Converts SetupResult enum to human-readable message.

\begin{lstlisting}[caption=getSetupResultMessage() Static Method]
static std::string getSetupResultMessage(SetupResult result);
\end{lstlisting}

\textbf{Parameters:}
\begin{itemize}
    \item \texttt{result} - SetupResult enum value
\end{itemize}

\textbf{Return Value:}
\begin{itemize}
    \item Human-readable error message string
\end{itemize}

\newpage
\section{Environment Variables}

\subsection{Overview}

The CGI handler sets environment variables according to RFC 3875. These variables provide the CGI script with information about the request, server, and client.

\subsection{Request Information Variables}

\begin{tabular}{|l|l|p{6cm}|}
\hline
\textbf{Variable} & \textbf{RFC Section} & \textbf{Description} \\
\hline
REQUEST\_METHOD & 4.1.12 & HTTP method (GET, POST, etc.) \\
\hline
REQUEST\_URI & Extension & Full request URI including query string \\
\hline
QUERY\_STRING & 4.1.7 & Query string portion of URI (after ?) \\
\hline
\end{tabular}

\subsection{Script Information Variables}

\begin{tabular}{|l|l|p{6cm}|}
\hline
\textbf{Variable} & \textbf{RFC Section} & \textbf{Description} \\
\hline
SCRIPT\_NAME & 4.1.13 & URI path to the script \\
\hline
SCRIPT\_FILENAME & Extension & Absolute filesystem path to script \\
\hline
PATH\_INFO & 4.1.5 & Extra path after script name \\
\hline
PATH\_TRANSLATED & 4.1.6 & DOCUMENT\_ROOT + PATH\_INFO \\
\hline
\end{tabular}

\textbf{Example:}

URI: \texttt{/cgi-bin/script.py/extra/path?query=value}

\begin{itemize}
    \item SCRIPT\_NAME = /cgi-bin/script.py
    \item PATH\_INFO = /extra/path
    \item QUERY\_STRING = query=value
\end{itemize}

\subsection{Server Information Variables}

\begin{tabular}{|l|l|p{6cm}|}
\hline
\textbf{Variable} & \textbf{RFC Section} & \textbf{Description} \\
\hline
SERVER\_NAME & 4.1.14 & Server hostname \\
\hline
SERVER\_PORT & 4.1.15 & Server port number \\
\hline
SERVER\_PROTOCOL & 4.1.16 & HTTP/1.0 \\
\hline
SERVER\_SOFTWARE & 4.1.17 & Webserv/1.0 \\
\hline
GATEWAY\_INTERFACE & 4.1.4 & CGI/1.1 \\
\hline
DOCUMENT\_ROOT & Extension & Web server document root \\
\hline
\end{tabular}

\textbf{Note:} DOCUMENT\_ROOT is currently set to the script directory. In production, this should be configured separately as the web server's document root.

\subsection{Content Variables}

\begin{tabular}{|l|l|p{6cm}|}
\hline
\textbf{Variable} & \textbf{RFC Section} & \textbf{Description} \\
\hline
CONTENT\_TYPE & 4.1.3 & MIME type of request body \\
\hline
CONTENT\_LENGTH & 4.1.2 & Length of request body in bytes \\
\hline
\end{tabular}

These are set from the Content-Type and Content-Length headers. If Content-Length is missing but a request body exists, it's calculated automatically.

\subsection{HTTP Header Variables}

All HTTP headers are converted to environment variables with the following transformation:

\begin{enumerate}
    \item Prefix with "HTTP\_"
    \item Convert to uppercase
    \item Replace hyphens with underscores
\end{enumerate}

\textbf{Examples:}
\begin{itemize}
    \item User-Agent $\rightarrow$ HTTP\_USER\_AGENT
    \item Accept-Language $\rightarrow$ HTTP\_ACCEPT\_LANGUAGE
    \item X-Custom-Header $\rightarrow$ HTTP\_X\_CUSTOM\_HEADER
\end{itemize}

\textbf{Multiple Headers:} Per RFC 3875, multiple headers with the same name are combined with commas.

Example: Two "Accept" headers are combined into HTTP\_ACCEPT = "text/html, application/json"

\subsection{Special Variables}

\begin{tabular}{|l|p{8cm}|}
\hline
\textbf{Variable} & \textbf{Description} \\
\hline
REDIRECT\_STATUS & Set to "200" for PHP compatibility \\
\hline
\end{tabular}

\subsection{Missing Variables (TODO)}

The following RFC 3875 variables are not yet implemented:

\begin{tabular}{|l|l|p{6cm}|}
\hline
\textbf{Variable} & \textbf{RFC Section} & \textbf{Reason} \\
\hline
AUTH\_TYPE & 4.1.1 & Authentication not implemented \\
\hline
REMOTE\_USER & 4.1.11 & Authentication not implemented \\
\hline
REMOTE\_IDENT & 4.1.10 & Authentication not implemented \\
\hline
REMOTE\_ADDR & 4.1.8 & Client info not available \\
\hline
REMOTE\_HOST & 4.1.9 & Client info not available \\
\hline
\end{tabular}

These will be added when authentication and client connection tracking are implemented in the web server.

\newpage
\section{Security Features}

\subsection{Overview}

The CGI implementation includes multiple security layers to prevent common vulnerabilities and attacks.

\subsection{Path Traversal Protection}

\subsubsection{Detection}

The \texttt{containsPathTraversal()} function detects:

\begin{itemize}
    \item \textbf{Directory traversal}: ".." sequences
    \item \textbf{Null byte injection}: Embedded null characters
\end{itemize}

\begin{lstlisting}[caption=Path Traversal Detection]
bool containsPathTraversal(const std::string& path) {
    // Check for null bytes
    if (path.find('\0') != std::string::npos) {
        return true;
    }
    
    // Check for ".." sequences
    if (path.find("..") != std::string::npos) {
        return true;
    }
    
    return false;
}
\end{lstlisting}

\subsubsection{Application}

Path traversal checks are applied to:
\begin{itemize}
    \item Script path (always checked in validateInputs())
    \item Interpreter path (checked if provided)
\end{itemize}

Any path containing ".." or null bytes is rejected with \texttt{SETUP\_ERROR\_PATH\_TRAVERSAL}.

\subsection{Input Sanitization}

\subsubsection{Environment Variable Sanitization}

All environment variable values are sanitized before being passed to the CGI script:

\begin{lstlisting}[caption=Environment Sanitization]
std::string sanitizeForEnv(const std::string& value) {
    std::string result;
    result.reserve(value.length());
    
    for (size_t i = 0; i < value.length(); ++i) {
        char c = value[i];
        // Remove dangerous characters
        if (c == '\0' || c == '\n' || c == '\r') {
            continue;
        }
        result += c;
    }
    
    return result;
}
\end{lstlisting}

\textbf{Removed characters:}
\begin{itemize}
    \item Null bytes (\textbackslash 0) - prevent string truncation attacks
    \item Newlines (\textbackslash n) - prevent header injection
    \item Carriage returns (\textbackslash r) - prevent header injection
\end{itemize}

\subsection{File System Security}

\subsubsection{Execute Permission Validation}

Scripts and interpreters must have execute permission:

\begin{lstlisting}[caption=Permission Check]
struct stat st;
if (stat(scriptPath.c_str(), &st) == 0) {
    if (!(st.st_mode & S_IXUSR)) {
        return SETUP_ERROR_SCRIPT_NOT_EXECUTABLE;
    }
}
\end{lstlisting}

\textbf{Note:} Only S\_IXUSR (user execute) is checked. This is sufficient for typical web server setups where scripts are owned by the web server user.

\subsubsection{Directory Change Restrictions}

Before executing scripts, the process changes to the script's directory. Root directory is excluded for security:

\begin{lstlisting}[caption=chdir Security]
std::string scriptDir = utils::getDirectory(_scriptPath);
if (!scriptDir.empty() && scriptDir != "/") {
    if (chdir(scriptDir.c_str()) == -1) {
        _exit(1);
    }
}
\end{lstlisting}

\textbf{Limitation:} Only root (/) is explicitly blocked. Other sensitive directories (/etc, /sys, /proc) should be protected via file system permissions.

\subsection{Port Validation}

Server port must be in valid range:

\begin{lstlisting}[caption=Port Validation]
bool isValidPort(int port) {
    return port >= 1 && port <= 65535;
}
\end{lstlisting}

Invalid ports are rejected with \texttt{SETUP\_ERROR\_INVALID\_PORT}.

\subsection{Overflow Protection}

\subsubsection{Integer Overflow}

Safe conversion functions prevent overflow:

\begin{lstlisting}[caption=Safe Integer Conversion]
int stringToInt(const std::string& str, int defaultValue) {
    std::istringstream iss(str);
    long value;
    
    if (!(iss >> value)) {
        return defaultValue;
    }
    
    // Check for overflow
    if (value > INT_MAX || value < INT_MIN) {
        return defaultValue;
    }
    
    return static_cast<int>(value);
}
\end{lstlisting}

\subsubsection{size\_t to String Conversion}

CONTENT\_LENGTH uses safe size\_t conversion to avoid casting overflow:

\begin{lstlisting}[caption=Safe size\_t Conversion]
std::string sizeToString(size_t value) {
    std::ostringstream oss;
    oss << value;
    return oss.str();
}
\end{lstlisting}

This prevents truncation when body size exceeds INT\_MAX (2GB).

\subsection{Process Isolation}

\subsubsection{Fork and Execute}

CGI scripts run in separate processes:
\begin{itemize}
    \item Isolated from web server process
    \item Cannot directly access server memory
    \item Can be killed without affecting server
\end{itemize}

\subsubsection{Error Handling}

Child process errors do not propagate to parent:
\begin{lstlisting}[caption=Child Process Error Handling]
if (pid == 0) {
    // Child process
    setupChildProcess(pipeToChild, pipeFromChild);
    _exit(1); // Only reached on error
}
\end{lstlisting}

Using \texttt{\_exit()} prevents cleanup functions from running in child, avoiding double-free bugs.

\subsection{Known Limitations}

\subsubsection{stderr Not Redirected}

Per RFC 3875 section 6.2, CGI script errors go to server's stderr. This is acceptable but may need enhancement for production systems requiring error capture.

\subsubsection{Directory Security}

Only root directory (/) is explicitly excluded from chdir. Sensitive directories should be protected via:
\begin{itemize}
    \item File system permissions
    \item AppArmor/SELinux policies
    \item chroot jails
\end{itemize}

\subsubsection{Resource Limits}

No built-in limits for:
\begin{itemize}
    \item CPU time (use ulimit or cgroups)
    \item Memory usage (use ulimit or cgroups)
    \item Number of processes (use ulimit)
\end{itemize}

These should be enforced at the system level.

\newpage
\section{Implementation Details}

\subsection{Execution Flow}

\subsubsection{High-Level Process}

\begin{enumerate}
    \item \textbf{Setup Phase}: Validate inputs and build environment
    \item \textbf{Execution Phase}: Fork, create pipes, execute script
    \item \textbf{Communication Phase}: Write request body, read output
    \item \textbf{Parsing Phase}: Parse CGI headers and body
    \item \textbf{Cleanup Phase}: Wait for child, close pipes
\end{enumerate}

\subsubsection{Detailed Flow Diagram}

\begin{verbatim}
Web Server                CGI Class              Child Process
    |                        |                         |
    |--[1] setup()---------->|                         |
    |                        |--validateInputs()       |
    |                        |--buildEnvironment()     |
    |<--[SETUP_OK]-----------| prepareEnvp()           |
    |                        |                         |
    |--[2] execute()-------->|                         |
    |                        |--createPipes()          |
    |                        |--fork()---------------->|
    |                        |                         |--chdir()
    |                        |                         |--dup2()
    |<--[pid, fds]-----------| (parent)                |--execve()
    |                        |                         |
    |--[3] write(body)------>|======================>  |
    |--close(writeFd)------->|                         |
    |                        |                         |
    |--[4] read(output)------|<=======================|
    |                        |                         |
    |--[5] parseCGIOutput()--|                         |
    |<--[headers,body]-------|                         |
    |                        |                         |
    |--[6] waitpid()---------|                         |--exit
    |                        |                         |
\end{verbatim}

\subsection{Pipe Management}

\subsubsection{Pipe Creation}

Two unidirectional pipes are created:

\begin{lstlisting}[caption=Pipe Creation]
int pipeToChild[2];    // Parent writes, child reads (stdin)
int pipeFromChild[2];  // Child writes, parent reads (stdout)

pipe(pipeToChild);
pipe(pipeFromChild);
\end{lstlisting}

\subsubsection{Parent Process}

Closes unused ends and sets non-blocking mode:

\begin{lstlisting}[caption=Parent Pipe Setup]
close(pipeToChild[0]);    // Close read end
close(pipeFromChild[1]);  // Close write end

// Set non-blocking (warnings logged, not critical)
fcntl(pipeToChild[1], F_SETFL, O_NONBLOCK);
fcntl(pipeFromChild[0], F_SETFL, O_NONBLOCK);

writePipeFd = pipeToChild[1];
readPipeFd = pipeFromChild[0];
\end{lstlisting}

\subsubsection{Child Process}

Redirects stdin/stdout and closes unused ends:

\begin{lstlisting}[caption=Child Pipe Setup]
close(pipeToChild[1]);    // Close write end
close(pipeFromChild[0]);  // Close read end

dup2(pipeToChild[0], STDIN_FILENO);   // Redirect stdin
dup2(pipeFromChild[1], STDOUT_FILENO); // Redirect stdout

close(pipeToChild[0]);
close(pipeFromChild[1]);
\end{lstlisting}

\subsection{Environment Preparation}

\subsubsection{Building the Map}

Environment variables are first stored in a \texttt{std::map<string, string>}:

\begin{lstlisting}[caption=Environment Map]
void addEnvironmentVariable(const std::string& key,
                           const std::string& value) {
    _env[key] = utils::sanitizeForEnv(value);
}
\end{lstlisting}

\subsubsection{Creating the Array}

The map is converted to a char* array for execve():

\begin{lstlisting}[caption=Environment Array Creation]
void prepareEnvp() {
    for (map::iterator it = _env.begin();
         it != _env.end(); ++it) {
        string envEntry = it->first + "=" + it->second;
        char* envStr = new char[envEntry.length() + 1];
        strcpy(envStr, envEntry.c_str());
        _envp.push_back(envStr);
    }
    _envp.push_back(NULL); // NULL terminator
}
\end{lstlisting}

\subsubsection{Memory Management}

Environment strings are cleaned up in destructor:

\begin{lstlisting}[caption=Environment Cleanup]
void cleanupEnvp() {
    // Skip NULL terminator
    for (size_t i = 0; i < _envp.size() && _envp[i] != NULL; ++i) {
        delete[] _envp[i];
    }
    _envp.clear();
}
\end{lstlisting}

\subsection{Script Execution}

\subsubsection{With Interpreter}

For scripts requiring an interpreter (Python, PHP, etc.):

\begin{lstlisting}[caption=Execution with Interpreter]
char* argv[3];
argv[0] = const_cast<char*>(_interpreterPath.c_str());
argv[1] = const_cast<char*>(scriptName.c_str());
argv[2] = NULL;

execve(_interpreterPath.c_str(), argv, _envp.data());
\end{lstlisting}

Example: \texttt{/usr/bin/python3 script.py}

\subsubsection{Direct Execution}

For executable scripts with shebang:

\begin{lstlisting}[caption=Direct Execution]
char* argv[2];
argv[0] = const_cast<char*>(scriptName.c_str());
argv[1] = NULL;

execve(scriptName.c_str(), argv, _envp.data());
\end{lstlisting}

\textbf{Critical:} scriptName must have proper lifetime extending until execve() call. Declaring it outside if/else blocks prevents undefined behavior.

\subsection{Header Parsing Algorithm}

\subsubsection{Header/Body Separation}

CGI output format per RFC 3875:

\begin{verbatim}
Header-Name: Value
Another-Header: Value

Body content starts here
\end{verbatim}

\subsubsection{Parsing Logic}

\begin{lstlisting}[caption=Header Parsing]
// Find separator (CRLF or LF)
size_t headerEnd = output.find("\r\n\r\n");
if (headerEnd == npos) {
    headerEnd = output.find("\n\n");
}

// Parse each header line
size_t pos = 0;
while (pos < headerEnd) {
    size_t lineEnd = output.find_first_of("\r\n", pos);
    string line = output.substr(pos, lineEnd - pos);
    
    size_t colonPos = line.find(':');
    if (colonPos != npos) {
        string key = utils::trim(line.substr(0, colonPos));
        string value = utils::trim(line.substr(colonPos + 1));
        headers[key] = value;
    }
    
    // Move to next line
    pos = lineEnd;
    if (output[pos] == '\r') pos++;
    if (output[pos] == '\n') pos++;
}

// Extract body
body = output.substr(headerEnd + separatorLen);
\end{lstlisting}

\textbf{Liberal Parsing:} Both key and value are trimmed to accept slightly malformed headers.

\subsection{URI Component Extraction}

\subsubsection{Script Name}

Everything before the query string:

\begin{lstlisting}[caption=Extract Script Name]
string extractScriptName(const string& uri) {
    size_t queryPos = uri.find('?');
    if (queryPos != npos) {
        return uri.substr(0, queryPos);
    }
    return uri;
}
\end{lstlisting}

\subsubsection{Query String}

Everything after the '?' character:

\begin{lstlisting}[caption=Extract Query String]
string extractQueryString(const string& uri) {
    size_t queryPos = uri.find('?');
    if (queryPos != npos && queryPos + 1 < uri.length()) {
        return uri.substr(queryPos + 1);
    }
    return "";
}
\end{lstlisting}

\subsubsection{Path Info}

Extra path after script name:

\begin{lstlisting}[caption=Extract Path Info]
string extractPathInfo(const string& uri,
                      const string& scriptName) {
    size_t scriptEnd = uri.find(scriptName);
    if (scriptEnd == npos) return "";
    
    size_t pathInfoStart = scriptEnd + scriptName.length();
    if (pathInfoStart >= uri.length()) return "";
    
    string pathInfo = uri.substr(pathInfoStart);
    
    // Remove query string if present
    size_t queryPos = pathInfo.find('?');
    if (queryPos != npos) {
        pathInfo = pathInfo.substr(0, queryPos);
    }
    
    return pathInfo;
}
\end{lstlisting}

\textbf{Note:} Uses find() which locates first occurrence. May be incorrect if script name appears multiple times in URI, but this is acceptable for typical usage.

\subsection{Non-Blocking I/O}

\subsubsection{Purpose}

Non-blocking pipes prevent deadlocks when:
\begin{itemize}
    \item CGI script output is large
    \item Request body is large
    \item Script reads and writes simultaneously
\end{itemize}

\subsubsection{Implementation}

\begin{lstlisting}[caption=Non-Blocking Configuration]
fcntl(fd, F_SETFL, O_NONBLOCK);
\end{lstlisting}

\textbf{Error Handling:} Errors are logged as warnings but not fatal. CGI works in blocking mode if fcntl() fails.

\subsubsection{Usage Pattern}

With non-blocking pipes, the caller must:
\begin{enumerate}
    \item Use select() or poll() to check readability/writability
    \item Handle EAGAIN/EWOULDBLOCK on read/write
    \item Retry operations when pipes become ready
\end{enumerate}

\newpage
\section{Testing Strategy}

\subsection{Test Suite Overview}

The CGI implementation includes 50 comprehensive tests located in \texttt{tests/test\_cgi.cpp}.

\subsection{Test Categories}

\subsubsection{libutils Tests (20 tests)}

\textbf{String Utilities:}
\begin{itemize}
    \item trim(), trimLeft(), trimRight() - whitespace handling
    \item toUpper(), toLower() - case conversion (3 tests)
    \item containsPathTraversal() - path security (4 tests including null byte)
\end{itemize}

\textbf{Path Utilities:}
\begin{itemize}
    \item getDirectory() - extract directory from path
    \item getFilename() - extract filename (5 tests)
    \item getExtension() - extract file extension
    \item hasExtension() - extension checking
\end{itemize}

\textbf{Validation:}
\begin{itemize}
    \item isValidPort() - port range validation
\end{itemize}

\textbf{Conversion:}
\begin{itemize}
    \item intToString() - int to string (5 tests including INT\_MAX/MIN)
    \item sizeToString() - size\_t to string (3 tests)
    \item stringToInt() - string to int with overflow checks (9 tests)
\end{itemize}

\subsubsection{libhttp Tests (4 tests)}

\begin{itemize}
    \item methodToString() - HTTP method conversion
    \item SetupResult messages - error message verification
\end{itemize}

\subsubsection{CGI Parsing Tests (6 tests)}

\begin{itemize}
    \item parseCGIOutput() - basic header/body separation
    \item Multiple headers - CRLF vs LF handling
    \item Empty output - edge case handling
    \item Missing headers - non-compliant CGI
    \item Header trimming - liberal parsing
    \item isCGIScript() - extension checking
\end{itemize}

\subsubsection{CGI Setup Tests (9 tests)}

\begin{itemize}
    \item Empty script path - validation
    \item Invalid port numbers - range checking
    \item Path traversal detection - security
    \item Non-existent script - file checking
    \item Non-executable script - permission checking
    \item Non-existent interpreter - validation
    \item Valid setup - success case
    \item Interpreter not found - new error code test
    \item Environment variable building - comprehensive check
\end{itemize}

\subsubsection{CGI Execution Tests (11 tests)}

\begin{itemize}
    \item Execute without setup - error handling
    \item Simple GET request - basic execution
    \item Simple POST request - body handling
    \item Query string parsing - QUERY\_STRING variable
    \item Environment variables - comprehensive check
    \item Multiple headers - HTTP\_* variable combination
    \item Path info extraction - PATH\_INFO and PATH\_TRANSLATED
    \item Content-Length calculation - automatic when missing
    \item Python CGI script - interpreter execution
    \item Shell CGI script - direct execution
    \item Non-blocking pipe mode - fcntl() configuration
\end{itemize}

\subsection{Test Execution}

\subsubsection{Building Tests}

\begin{lstlisting}[language=bash, caption=Build and Run Tests]
cd /home/mkurkar/QuantumBytes/web-server/
make fclean && make test
\end{lstlisting}

\subsubsection{Test Output Format}

\begin{verbatim}
Running: test_name... PASSED
Running: test_name... FAILED
  FAILED: condition is false

================================
Test Results:
  Total: 50
  Passed: 50
  Failed: 0
================================
\end{verbatim}

\subsection{Test CGI Scripts}

Sample scripts in \texttt{tests/cgi-bin/}:

\subsubsection{hello.py}

Simple Python CGI:

\begin{lstlisting}[language=Python, caption=hello.py]
#!/usr/bin/env python3
print("Content-Type: text/plain")
print()
print("Hello from Python CGI!")
\end{lstlisting}

\subsubsection{echo\_env.py}

Environment variable display:

\begin{lstlisting}[language=Python, caption=echo\_env.py]
#!/usr/bin/env python3
import os
print("Content-Type: text/plain")
print()
for key, value in sorted(os.environ.items()):
    print(f"{key}={value}")
\end{lstlisting}

\subsubsection{post\_echo.py}

POST body echo:

\begin{lstlisting}[language=Python, caption=post\_echo.py]
#!/usr/bin/env python3
import sys
print("Content-Type: text/plain")
print()
body = sys.stdin.read()
print(f"Received: {body}")
\end{lstlisting}

\subsubsection{simple.sh}

Shell script CGI:

\begin{lstlisting}[language=bash, caption=simple.sh]
#!/bin/bash
echo "Content-Type: text/plain"
echo ""
echo "Hello from shell CGI!"
echo "QUERY_STRING=$QUERY_STRING"
\end{lstlisting}

\subsection{Coverage Analysis}

\subsubsection{Code Coverage}

All public methods tested:
\begin{itemize}
    \item setup() - 100\% path coverage (all 7 SetupResult codes)
    \item execute() - success and error paths
    \item parseCGIOutput() - various formats and edge cases
    \item isCGIScript() - positive and negative cases
    \item getError() - error retrieval
    \item getSetupResultMessage() - all enum values
\end{itemize}

\subsubsection{Edge Cases}

\begin{itemize}
    \item Empty strings
    \item Null bytes in paths
    \item Integer overflow (INT\_MAX, INT\_MIN)
    \item size\_t overflow scenarios
    \item Missing headers
    \item Malformed CGI output
    \item Non-blocking vs blocking mode
    \item Multiple identical headers
\end{itemize}

\subsubsection{Security Coverage}

\begin{itemize}
    \item Path traversal (..)
    \item Null byte injection
    \item Port validation
    \item Execute permission checking
    \item Environment sanitization
    \item Header injection prevention
\end{itemize}

\subsection{Test Assertions}

\begin{lstlisting}[caption=Test Macros]
ASSERT_TRUE(condition)
ASSERT_FALSE(condition)
ASSERT_EQ(expected, actual)
ASSERT_NE(expected, actual)
ASSERT_CONTAINS(haystack, needle)
\end{lstlisting}

All assertions throw on failure with colored output for easy debugging.

\newpage
\section{Usage Examples}

\subsection{Basic GET Request}

\begin{lstlisting}[caption=Simple GET Request Handler]
#include "CGI.hpp"
#include "Request.hpp"
#include <unistd.h>
#include <sys/wait.h>

void handleCGIRequest(const libhttp::Request& request) {
    CGI cgi;
    
    // Setup CGI
    CGI::SetupResult result = cgi.setup(
        request,
        "/var/www/cgi-bin/script.py",  // Script path
        "/usr/bin/python3",              // Interpreter
        "example.com",                   // Server name
        8080                             // Server port
    );
    
    if (result != CGI::SETUP_OK) {
        std::cerr << "Setup failed: "
                  << CGI::getSetupResultMessage(result)
                  << std::endl;
        return;
    }
    
    // Execute CGI
    int readFd, writeFd;
    pid_t pid = cgi.execute(readFd, writeFd);
    
    if (pid == -1) {
        std::cerr << "Execute failed: "
                  << cgi.getError() << std::endl;
        return;
    }
    
    // Close write pipe immediately for GET
    close(writeFd);
    
    // Read CGI output
    std::string output;
    char buffer[4096];
    ssize_t n;
    
    while ((n = read(readFd, buffer, sizeof(buffer))) > 0) {
        output.append(buffer, n);
    }
    
    close(readFd);
    
    // Wait for child process
    int status;
    waitpid(pid, &status, 0);
    
    // Parse output
    std::map<std::string, std::string> headers;
    std::string body;
    
    if (CGI::parseCGIOutput(output, headers, body)) {
        // Send headers to client
        for (map::iterator it = headers.begin();
             it != headers.end(); ++it) {
            sendHeader(it->first, it->second);
        }
        // Send body
        sendBody(body);
    }
}
\end{lstlisting}

\subsection{POST Request with Body}

\begin{lstlisting}[caption=POST Request Handler]
void handlePOSTRequest(const libhttp::Request& request) {
    CGI cgi;
    
    // Setup
    CGI::SetupResult result = cgi.setup(
        request,
        "/var/www/cgi-bin/form_handler.py",
        "/usr/bin/python3",
        "example.com",
        8080
    );
    
    if (result != CGI::SETUP_OK) {
        handleError(result);
        return;
    }
    
    // Execute
    int readFd, writeFd;
    pid_t pid = cgi.execute(readFd, writeFd);
    
    if (pid == -1) {
        handleError(cgi.getError());
        return;
    }
    
    // Write request body to CGI
    const std::string& body = request.getBody();
    if (!body.empty()) {
        ssize_t written = write(writeFd, body.c_str(),
                               body.length());
        if (written == -1) {
            std::cerr << "Failed to write body" << std::endl;
        }
    }
    
    // CRITICAL: Close write pipe to signal EOF
    close(writeFd);
    
    // Read output (same as GET example)
    std::string output = readAllFromFd(readFd);
    close(readFd);
    
    // Wait and parse
    waitpid(pid, NULL, 0);
    processOutput(output);
}
\end{lstlisting}

\subsection{Python CGI Script Example}

\begin{lstlisting}[language=Python, caption=Form Processing Script]
#!/usr/bin/env python3
import sys
import os
import urllib.parse

# Read POST body
content_length = int(os.environ.get('CONTENT_LENGTH', 0))
post_data = sys.stdin.read(content_length) if content_length > 0 else ''

# Parse form data
params = urllib.parse.parse_qs(post_data)

# Generate response
print("Content-Type: text/html")
print()
print("<html><body>")
print("<h1>Form Submission Result</h1>")

for key, values in params.items():
    print(f"<p>{key}: {', '.join(values)}</p>")

print("</body></html>")
\end{lstlisting}

\subsection{Shell Script Example}

\begin{lstlisting}[language=bash, caption=Environment Display Script]
#!/bin/bash

echo "Content-Type: text/plain"
echo ""
echo "CGI Environment Variables"
echo "========================="
echo ""
echo "Request Method: $REQUEST_METHOD"
echo "Request URI: $REQUEST_URI"
echo "Query String: $QUERY_STRING"
echo "Script Name: $SCRIPT_NAME"
echo "Path Info: $PATH_INFO"
echo "Server Name: $SERVER_NAME"
echo "Server Port: $SERVER_PORT"
echo "Content Type: $CONTENT_TYPE"
echo "Content Length: $CONTENT_LENGTH"
\end{lstlisting}

\subsection{Error Handling Example}

\begin{lstlisting}[caption=Comprehensive Error Handling]
bool executeCGI(const libhttp::Request& request,
                const std::string& scriptPath) {
    CGI cgi;
    
    // Determine interpreter from extension
    std::string interpreter;
    if (CGI::isCGIScript(scriptPath, ".py")) {
        interpreter = "/usr/bin/python3";
    } else if (CGI::isCGIScript(scriptPath, ".pl")) {
        interpreter = "/usr/bin/perl";
    } else if (CGI::isCGIScript(scriptPath, ".php")) {
        interpreter = "/usr/bin/php-cgi";
    }
    // else: direct execution (script must have shebang)
    
    // Setup with detailed error checking
    CGI::SetupResult result = cgi.setup(
        request, scriptPath, interpreter,
        getServerName(), getServerPort()
    );
    
    switch (result) {
        case CGI::SETUP_OK:
            break;
            
        case CGI::SETUP_ERROR_EMPTY_SCRIPT_PATH:
            sendError(500, "Internal error: empty script path");
            return false;
            
        case CGI::SETUP_ERROR_INVALID_PORT:
            sendError(500, "Internal error: invalid port");
            return false;
            
        case CGI::SETUP_ERROR_PATH_TRAVERSAL:
            sendError(400, "Bad request: invalid script path");
            return false;
            
        case CGI::SETUP_ERROR_SCRIPT_NOT_FOUND:
            sendError(404, "Script not found");
            return false;
            
        case CGI::SETUP_ERROR_SCRIPT_NOT_EXECUTABLE:
            sendError(500, "Script not executable");
            return false;
            
        case CGI::SETUP_ERROR_INTERPRETER_NOT_FOUND:
            sendError(500, "Interpreter not found");
            return false;
            
        default:
            sendError(500, "Unknown setup error");
            return false;
    }
    
    // Execute
    int readFd, writeFd;
    pid_t pid = cgi.execute(readFd, writeFd);
    
    if (pid == -1) {
        sendError(500, "Failed to execute CGI: " + cgi.getError());
        return false;
    }
    
    // Process request (write body, read output, etc.)
    bool success = processCGIRequest(pid, readFd, writeFd, request);
    
    return success;
}
\end{lstlisting}

\subsection{Non-Blocking I/O Example}

\begin{lstlisting}[caption=Non-Blocking CGI Communication]
#include <sys/select.h>

bool readCGIOutput(int readFd, std::string& output,
                   int timeoutMs = 30000) {
    fd_set readSet;
    struct timeval timeout;
    
    timeout.tv_sec = timeoutMs / 1000;
    timeout.tv_usec = (timeoutMs % 1000) * 1000;
    
    while (true) {
        FD_ZERO(&readSet);
        FD_SET(readFd, &readSet);
        
        int ret = select(readFd + 1, &readSet, NULL, NULL,
                        &timeout);
        
        if (ret == -1) {
            return false; // Error
        } else if (ret == 0) {
            return false; // Timeout
        }
        
        if (FD_ISSET(readFd, &readSet)) {
            char buffer[4096];
            ssize_t n = read(readFd, buffer, sizeof(buffer));
            
            if (n == -1) {
                if (errno == EAGAIN || errno == EWOULDBLOCK) {
                    continue; // Try again
                }
                return false; // Real error
            } else if (n == 0) {
                return true; // EOF - success
            }
            
            output.append(buffer, n);
        }
    }
}
\end{lstlisting}

\newpage
\section{Integration Guide}

\subsection{Project Structure}

\begin{verbatim}
web-server/
|------ Makefile                          # Main build file
|------ dependencies/
|     |------ libutils/                     # General utilities
|     |     |------ include/StringUtils.hpp
|     |     |------ src/Utils.cpp
|     |     `------ Makefile
|     `------ libhttp/                      # HTTP library
|         |------ include/
|         |     |------ CGI.hpp
|         |     |------ Request.hpp
|         |     |------ Response.hpp
|         |     |------ Utils.hpp
|         |     `------ Constants.hpp
|         |------ src/
|         |     |------ CGI.cpp
|         |     |------ Request.cpp
|         |     |------ Response.cpp
|         |     `------ Utils.cpp
|         `------ Makefile
`------ tests/
    |------ test_cgi.cpp
    `------ cgi-bin/
        |------ hello.py
        |------ echo_env.py
        |------ post_echo.py
        `------ simple.sh
\end{verbatim}

\subsection{Makefile Configuration}

\subsubsection{Main Makefile}

\begin{lstlisting}[language=make, caption=Top-Level Makefile]
NAME = webserv
CXX = c++
CXXFLAGS = -Wall -Wextra -Werror -std=c++98

# Dependencies (order matters!)
DEPENDENCIES = libhttp libutils

# Include paths for dependencies
INCLUDES = -I dependencies/libhttp/include \
           -I dependencies/libutils/include

# Library paths
LIBDIRS = -L dependencies/libhttp \
          -L dependencies/libutils

# Link libraries
LIBS = -lhttp -lutils

# Build dependencies first
$(NAME): $(DEPENDENCIES)
	$(CXX) $(CXXFLAGS) $(INCLUDES) \
	       src/*.cpp $(LIBDIRS) $(LIBS) -o $(NAME)

libhttp:
	$(MAKE) -C dependencies/libhttp

libutils:
	$(MAKE) -C dependencies/libutils
\end{lstlisting}

\textbf{Important:} Library linking order matters. libhttp depends on libutils, so link as: \texttt{-lhttp -lutils}

\subsubsection{Library Makefiles}

Each library has its own Makefile building a static library (.a file).

\subsection{Build Instructions}

\subsubsection{Building the Project}

\begin{lstlisting}[language=bash, caption=Build Commands]
# Clean and build everything
make fclean && make

# Build only
make

# Clean
make clean    # Remove objects
make fclean   # Remove objects and binaries

# Build and run tests
make test
\end{lstlisting}

\subsubsection{Build Dependencies}

Required:
\begin{itemize}
    \item C++ compiler with C++98 support (g++, clang++)
    \item GNU Make
    \item Standard C++ library
\end{itemize}

Optional (for tests):
\begin{itemize}
    \item Python 3 (for Python CGI scripts)
    \item Bash (for shell CGI scripts)
\end{itemize}

\subsection{Integration Steps}

\subsubsection{Step 1: Include Headers}

\begin{lstlisting}[caption=Include CGI Headers]
#include "CGI.hpp"
#include "Request.hpp"
\end{lstlisting}

\subsubsection{Step 2: Parse HTTP Request}

Use the Request class to parse incoming HTTP requests:

\begin{lstlisting}[caption=Request Parsing]
libhttp::Request request;
// Parse raw HTTP request into request object
// (implementation depends on your HTTP parser)
\end{lstlisting}

\subsubsection{Step 3: Detect CGI Requests}

Check if the requested resource is a CGI script:

\begin{lstlisting}[caption=CGI Detection]
std::string uri = request.getUri();
std::string scriptPath = mapURIToFilePath(uri);

bool isCGI = CGI::isCGIScript(scriptPath, ".py") ||
             CGI::isCGIScript(scriptPath, ".pl") ||
             CGI::isCGIScript(scriptPath, ".php") ||
             CGI::isCGIScript(scriptPath, ".cgi");

if (isCGI) {
    handleCGIRequest(request, scriptPath);
} else {
    handleStaticFile(scriptPath);
}
\end{lstlisting}

\subsubsection{Step 4: Execute CGI}

\begin{lstlisting}[caption=CGI Execution]
CGI cgi;

// Setup
CGI::SetupResult result = cgi.setup(
    request,
    scriptPath,
    getInterpreterForScript(scriptPath),
    serverConfig.hostname,
    serverConfig.port
);

if (result != CGI::SETUP_OK) {
    handleSetupError(result);
    return;
}

// Execute
int readFd, writeFd;
pid_t pid = cgi.execute(readFd, writeFd);

if (pid == -1) {
    handleExecuteError(cgi.getError());
    return;
}

// Communicate with CGI process
processCGI(pid, readFd, writeFd, request);
\end{lstlisting}

\subsubsection{Step 5: Process CGI Output}

\begin{lstlisting}[caption=Output Processing]
// Read output
std::string output = readAllFromPipe(readFd);

// Parse headers and body
std::map<std::string, std::string> headers;
std::string body;

if (CGI::parseCGIOutput(output, headers, body)) {
    // Build HTTP response
    sendHTTPHeaders(headers);
    sendHTTPBody(body);
} else {
    sendError(500, "CGI output parsing failed");
}

// Clean up
close(readFd);
waitpid(pid, NULL, 0);
\end{lstlisting}

\subsection{Configuration Recommendations}

\subsubsection{CGI Directory Structure}

\begin{verbatim}
/var/www/
├---- html/              # Static files
`------ cgi-bin/          # CGI scripts
    ├---- scripts/
    |     ├---- form.py
    |     `------ api.pl
    `------ interpreters/
        `------ custom_interpreter
\end{verbatim}

\subsubsection{File Permissions}

\begin{lstlisting}[language=bash, caption=Permission Setup]
# Make scripts executable
chmod +x /var/www/cgi-bin/*.py
chmod +x /var/www/cgi-bin/*.pl
chmod +x /var/www/cgi-bin/*.sh

# Secure permissions (owner execute only)
chmod 700 /var/www/cgi-bin/scripts/

# Set ownership to web server user
chown -R www-data:www-data /var/www/cgi-bin/
\end{lstlisting}

\subsubsection{Interpreter Paths}

Common interpreter locations:

\begin{tabular}{|l|l|}
\hline
\textbf{Language} & \textbf{Path} \\
\hline
Python 3 & /usr/bin/python3 \\
Perl & /usr/bin/perl \\
PHP & /usr/bin/php-cgi \\
Ruby & /usr/bin/ruby \\
\hline
\end{tabular}

Verify interpreter paths:
\begin{lstlisting}[language=bash]
which python3
which perl
which php-cgi
\end{lstlisting}

\newpage
\section{Troubleshooting}

\subsection{Common Issues}

\subsubsection{CGI Script Hangs}

\textbf{Symptom:} CGI process never completes, server hangs waiting

\textbf{Cause:} Write pipe not closed after sending request body

\textbf{Solution:}
\begin{lstlisting}
// Always close writePipeFd after writing
write(writePipeFd, body.c_str(), body.length());
close(writePipeFd);  // CRITICAL - signals EOF to CGI
\end{lstlisting}

For GET requests, close immediately after execute():
\begin{lstlisting}
pid_t pid = cgi.execute(readFd, writeFd);
close(writeFd);  // Close immediately for GET
\end{lstlisting}

\subsubsection{Script Not Found Error}

\textbf{Symptom:} \texttt{SETUP\_ERROR\_SCRIPT\_NOT\_FOUND}

\textbf{Causes:}
\begin{itemize}
    \item Script path is incorrect
    \item Script file doesn't exist
    \item Insufficient permissions to stat() the file
\end{itemize}

\textbf{Solutions:}
\begin{lstlisting}[language=bash]
# Check file exists
ls -l /path/to/script.py

# Check permissions
chmod +r /path/to/script.py

# Verify absolute path
readlink -f /path/to/script.py
\end{lstlisting}

\subsubsection{Script Not Executable}

\textbf{Symptom:} \texttt{SETUP\_ERROR\_SCRIPT\_NOT\_EXECUTABLE}

\textbf{Cause:} Script lacks execute permission

\textbf{Solution:}
\begin{lstlisting}[language=bash]
chmod +x /path/to/script.py
\end{lstlisting}

For web server user (www-data):
\begin{lstlisting}[language=bash]
chown www-data:www-data /path/to/script.py
chmod 700 /path/to/script.py
\end{lstlisting}

\subsubsection{Interpreter Not Found}

\textbf{Symptom:} \texttt{SETUP\_ERROR\_INTERPRETER\_NOT\_FOUND}

\textbf{Causes:}
\begin{itemize}
    \item Interpreter path is wrong
    \item Interpreter not installed
    \item Interpreter not executable
\end{itemize}

\textbf{Solutions:}
\begin{lstlisting}[language=bash]
# Find interpreter
which python3
which perl

# Install if missing
apt-get install python3
apt-get install perl

# Check permissions
ls -l /usr/bin/python3
\end{lstlisting}

\subsubsection{Path Traversal Error}

\textbf{Symptom:} \texttt{SETUP\_ERROR\_PATH\_TRAVERSAL}

\textbf{Cause:} Script path contains ".." or null bytes

\textbf{Solution:} Sanitize user input before passing to CGI:
\begin{lstlisting}
std::string sanitizePath(const std::string& path) {
    // Remove ".." sequences
    // Validate against whitelist
    // Use realpath() to resolve symlinks
}
\end{lstlisting}

\subsubsection{Empty CGI Output}

\textbf{Symptom:} No output received from CGI script

\textbf{Causes:}
\begin{itemize}
    \item Script has syntax error
    \item Script exits immediately
    \item Script output goes to stderr
    \item Timeout too short
\end{itemize}

\textbf{Debug:}
\begin{lstlisting}[language=bash]
# Test script manually
cd /path/to/script/dir
./script.py

# Check for errors
python3 -m py_compile script.py

# Check shebang
head -n 1 script.py
\end{lstlisting}

\subsubsection{Parse Error}

\textbf{Symptom:} \texttt{parseCGIOutput()} returns false

\textbf{Cause:} CGI output missing header/body separator

\textbf{Solution:} Ensure CGI script outputs blank line:
\begin{lstlisting}[language=Python]
print("Content-Type: text/plain")
print()  # CRITICAL - blank line separates headers from body
print("Body content")
\end{lstlisting}

\subsection{Debugging Techniques}

\subsubsection{Enable Debug Logging}

Add logging to CGI execution:

\begin{lstlisting}
std::cerr << "Setup result: " << result << std::endl;
std::cerr << "PID: " << pid << std::endl;
std::cerr << "Output length: " << output.length() << std::endl;
\end{lstlisting}

\subsubsection{Test Scripts Manually}

Run scripts outside web server:

\begin{lstlisting}[language=bash]
# Set environment manually
export REQUEST_METHOD=GET
export QUERY_STRING="key=value"
export SCRIPT_NAME="/cgi-bin/script.py"

# Run script
cd /var/www/cgi-bin/
./script.py

# For POST
echo "name=value" | ./script.py
\end{lstlisting}

\subsubsection{Check Child Process Status}

\begin{lstlisting}
int status;
waitpid(pid, &status, 0);

if (WIFEXITED(status)) {
    int exitCode = WEXITSTATUS(status);
    std::cerr << "Exit code: " << exitCode << std::endl;
}

if (WIFSIGNALED(status)) {
    int signal = WTERMSIG(status);
    std::cerr << "Killed by signal: " << signal << std::endl;
}
\end{lstlisting}

\subsubsection{Capture stderr}

Redirect CGI stderr to a file:

\begin{lstlisting}
// In child process, before execve()
int errFd = open("/tmp/cgi_errors.log",
                 O_WRONLY | O_CREAT | O_APPEND, 0644);
dup2(errFd, STDERR_FILENO);
close(errFd);
\end{lstlisting}

\subsubsection{Use strace}

Trace system calls:

\begin{lstlisting}[language=bash]
# Trace web server process
strace -f -o trace.log ./webserv

# Filter CGI-related calls
grep -E "fork|execve|pipe|read|write" trace.log
\end{lstlisting}

\subsection{Performance Issues}

\subsubsection{Slow CGI Execution}

\textbf{Causes:}
\begin{itemize}
    \item Script is slow (algorithmic complexity)
    \item Large output buffering
    \item Database queries
    \item External API calls
\end{itemize}

\textbf{Solutions:}
\begin{itemize}
    \item Profile script performance
    \item Add caching
    \item Use faster interpreter
    \item Implement timeout mechanism
\end{itemize}

\subsubsection{Memory Leaks}

\textbf{Check with Valgrind:}

\begin{lstlisting}[language=bash]
valgrind --leak-check=full \
         --show-leak-kinds=all \
         --track-fds=yes \
         ./webserv
\end{lstlisting}

Common leaks:
\begin{itemize}
    \item Environment strings not deleted (cleanupEnvp())
    \item Pipe file descriptors not closed
    \item Child processes not reaped (waitpid())
\end{itemize}

\subsection{Error Messages Reference}

\begin{tabular}{|p{6cm}|p{8cm}|}
\hline
\textbf{Error Message} & \textbf{Meaning \& Solution} \\
\hline
Script path cannot be empty &
Script path parameter is empty. Provide valid path. \\
\hline
Invalid port number &
Port outside 1-65535 range. Check server configuration. \\
\hline
Path traversal detected &
Path contains ".." or null bytes. Sanitize user input. \\
\hline
Script file not found &
File doesn't exist. Check path and permissions. \\
\hline
Script file is not executable &
Missing execute permission. Use chmod +x. \\
\hline
Interpreter not found &
Interpreter path wrong or not installed. Check with which. \\
\hline
Failed to fork &
System out of resources. Check ulimit -u. \\
\hline
Failed to create pipe &
System resource limit. Check ulimit -n. \\
\hline
Warning: failed to set non-blocking &
fcntl() error. CGI works but in blocking mode. \\
\hline
\end{tabular}

\subsection{Best Practices}

\begin{itemize}
    \item Always close writePipeFd after writing
    \item Always waitpid() to reap child processes
    \item Set reasonable timeouts (30s typical)
    \item Validate all user input before CGI
    \item Use absolute paths for scripts and interpreters
    \item Test scripts manually before deploying
    \item Monitor system resource limits
    \item Implement proper error handling
    \item Log CGI errors for debugging
    \item Use non-blocking I/O with select()/poll()
\end{itemize}

\newpage
\section{Conclusion}

This CGI implementation provides a robust, RFC 3875 compliant solution for executing external scripts in a web server environment. Key achievements:

\begin{itemize}
    \item \textbf{Comprehensive}: All required CGI meta-variables implemented
    \item \textbf{Secure}: Multiple security layers (path traversal, sanitization, permissions)
    \item \textbf{Tested}: 50 comprehensive tests covering all functionality
    \item \textbf{Well-Documented}: Clear code comments and this documentation
    \item \textbf{Maintainable}: Clean separation of concerns (SRP, DRY, KISS)
    \item \textbf{C++98 Compliant}: Meets 42 School requirements
\end{itemize}

\subsection{Future Enhancements}

Potential improvements for production use:

\begin{itemize}
    \item Authentication support (AUTH\_TYPE, REMOTE\_USER)
    \item Client connection tracking (REMOTE\_ADDR, REMOTE\_HOST)
    \item stderr capture and logging
    \item Resource limits (CPU, memory, timeout)
    \item FastCGI protocol support
    \item Connection pooling for CGI processes
    \item Async I/O with event loops (epoll, kqueue)
\end{itemize}

\subsection{References}

\begin{itemize}
    \item RFC 3875 - The Common Gateway Interface (CGI) Version 1.1
    \item POSIX.1-2001 - fork(), execve(), pipe(), dup2()
    \item C++98 Standard (ISO/IEC 14882:1998)
    \item 42 School Webserv Project Subject
\end{itemize}

\end{document}
